
Token log-probabilities --- computed without additional cost during generation by any autoregressive LLM --- are widely used as zero-cost hallucination detection signals, yet the aggregation step that converts a per-token log-probability sequence into a single confidence score has not been studied systematically. This paper demonstrates that this choice is not arbitrary: the optimal aggregation function is determined by hallucination type, and choosing the wrong function costs up to 12 AUROC points. Through a controlled ablation isolating aggregation function as the sole experimental variable, three-step causal mechanism is established --- hallucination type determines token probability distribution shape (peaked for recall-failure, flat for imitative-falsehood), which determines aggregation function optimality (minimum vs. mean) --- confirmed by distributional, rank-correlation, and threshold-level evidence across LLaMA-2-7B and Mistral-7B-v0.1. On TriviaQA, minimum log-probability outperforms mean by 5.6--11.9 AUROC points; on TruthfulQA, mean outperforms minimum by 11.1--12.0 AUROC points --- both with bootstrap 95\% confidence intervals excluding zero. Additionally, raw unnormalized log-probability sum achieves the highest AUROC on TriviaQA (0.895--0.896), exceeding published Semantic Entropy AUROC ($\approx 0.79)$ at one-tenth the inference cost under this evaluation protocol. These findings provide the first mechanism-grounded aggregation selection criterion for zero-cost hallucination detection.

